\documentclass{article}
\usepackage{mathtools} 
\usepackage{fontspec}
\usepackage[UTF8]{ctex}
\usepackage{amsthm}
\usepackage{mdframed}
\usepackage{xcolor}
\usepackage{amssymb}
\usepackage{amsmath}


% 定义新的带灰色背景的说明环境 zremark
\newmdtheoremenv[
  backgroundcolor=gray!10,
  % 边框与背景一致，边框线会消失
  linecolor=gray!10
]{zremark}{说明}

% 通用矩阵命令: \flexmatrix{矩阵名}{元素符号}{行数}{列数}
\newcommand{\flexmatrix}[4]{
  \[
  #1 = \begin{pmatrix}
    #2_{11}     & #2_{12}     & \cdots & #2_{1#4}   \\
    #2_{21}     & #2_{22}     & \cdots & #2_{2#4}   \\
    \vdots      & \vdots      & \ddots & \vdots     \\
    #2_{#31}    & #2_{#32}    & \cdots & #2_{#3#4}
  \end{pmatrix}
  \]
}

% 简化版命令（默认矩阵名为A，元素符号为a）: \quickmatrix{行数}{列数}
\newcommand{\quickmatrix}[2]{\flexmatrix{A}{a}{#1}{#2}}

\begin{document}
\title{习题 9.5}
\author{张志聪}
\maketitle

\section*{1}

f为连续函数，由微积分学基本定理可知，对任意一点$a$，我们有
\begin{align*}
  F(x) = \int_a^x f(t) dt
\end{align*}
于是
\begin{align*}
  \frac{d}{dx} \int_{u(x)}^{v(x)} f(t) dt
   & = \frac{d}{dx} \left(\int_{u(x)}^{a} f(t) dt + \frac{d}{dx} \int_{a}^{v(x)} f(t) dt\right)   \\
   & = \frac{d}{dx} \left(- \int_{a}^{u(x)} f(t) dt + \frac{d}{dx} \int_{a}^{v(x)} f(t) dt\right) \\
   & = \frac{d}{dx} \left( - F(u(x)) + F(v(x)) \right)                                            \\
   & = - F'(u(x))u'(x) + F'(v(x))v'(x)
\end{align*}
严格的来说$\int_{a}^{u(x)} f(t) dt = F(u(x))$可能需要证明，
可以利用定理9.9，定理9.10的证明过程，来证明把积分上限改为$u(x)$也是成立的，
过程比较多，这里省略。

\section*{2}
\begin{align*}
  F(x) & = \int_a^x f(t) (x - t) dt               \\
       & = \int_a^x f(t)x - f(t)t dt              \\
       & = \int_a^x f(t)x dt - \int_a^x f(t)t dt  \\
       & = x \int_a^x f(t) dt - \int_a^x f(t)t dt
\end{align*}
利用微积分学基本定理
\begin{align*}
  F'(x) & = \int_a^x f(t) dt + x f(x) - f(x) x \\
        & = \int_a^x f(t) dt
\end{align*}
于是
\begin{align*}
  F''(x) = f(x)
\end{align*}

\section*{3}

\begin{itemize}
  \item (1)

        \begin{itemize}
          \item 方法一：利用微积分学基本定理

                因为$cos t^2$函数在$\mathbb{R}$上连续，
                利用积分第一中值定理，存在$\xi \in (0, x)$，使得
                \begin{align*}
                  \int_0^x cos t^2 dt = cos \xi^2 (x - 0) = x cos \xi^2
                \end{align*}
                所以
                \begin{align*}
                  \lim \limits_{x \to 0} \frac{1}{x} \int_0^x cos t^2 dt
                   & = \lim \limits_{x \to 0} \frac{1}{x} x cos \xi^2 \\
                   & = \lim \limits_{x \to 0} cos \xi^2
                \end{align*}
                因为$x \to 0$，$\xi \to 0$，所以
                \begin{align*}
                  \lim \limits_{x \to 0} cos \xi^2
                   & = cos 0 \\
                   & = 1
                \end{align*}

          \item 方法二：利用微积分学基本定理

                因为$cos t^2$函数在$\mathbb{R}$上连续，
                由微积分学基本定理，$cos t^2$有原函数，设为$F(t)$，
                又由定理9.9可知，$F(x)$是连续的，$\lim\limits_{x \to 0} F(x) = F(0)$，
                于是利用洛必达定义
                \begin{align*}
                  \lim \limits_{x \to 0} \frac{1}{x} \int_0^x cos t^2 dt
                   & = \lim \limits_{x \to 0} \frac{F(x) - F(0)}{x} \\
                   & = \lim \limits_{x \to 0} \frac{cos x^2}{1}     \\
                   & = 1
                \end{align*}
        \end{itemize}

  \item (2)

        因为$e^x \geq 1, x \in [0, + \infty)$，由积分不等式可知，
        分式是未定式$\frac{\infty}{\infty}$，使用洛必达法则
        \begin{align*}
          \lim \limits_{x \to 0} \frac{\left(\int_0^x e^{t^2} dt\right)^2}{\int_0^x e^{2t^2} dt}
           & = \lim \limits_{x \to 0} \frac{2 e^{x^2}\int_0^x e^{t^2} dt}{e^{2x^2}} \\
           & = \lim \limits_{x \to 0} \frac{2 \int_0^x e^{t^2} dt}{e^{x^2}}         \\
           & = \lim \limits_{x \to 0} \frac{2 e^{x^2}}{2x e^{x^2}}                  \\
           & = \lim \limits_{x \to 0} \frac{1}{x}                                   \\
           & = 0
        \end{align*}

\end{itemize}

\section*{4}

\begin{itemize}
  \item (2)

        提示：$x = 2sinx$.
\end{itemize}

\section*{5}

\begin{itemize}
  \item (1)
        \begin{align*}
          \int_{-a}^{a} f(x) dx
           & = \int_{-a}^{0} f(x) dx + \int_{0}^{a} f(x) dx
        \end{align*}
        对第一个定积分使用换元积分法，令$x = -t, dx = -dt$，
        当$t$由$a$变为$0$时，$x$由$-a$增到$0$，于是
        \begin{align*}
          \int_{-a}^{0} f(x) dx & = - \int_{a}^{0} f(-t) dt \\
                                & = \int_{0}^{a} f(-t) dt   \\
                                & = \int_{0}^{a} f(-x) dx
        \end{align*}
        由奇函数可知$f(x) + f(-x) = 0$，所以
        \begin{align*}
          \int_{-a}^{a} f(x) dx
           & = \int_{-a}^{0} f(x) dx + \int_{0}^{a} f(x) dx \\
           & = \int_{0}^{a} f(-x) dx + \int_{0}^{a} f(x) dx \\
           & = \int_{0}^{a} f(-x) + f(x) dx                 \\
           & = \int_{0}^{a} 0 dx                            \\
           & = 0
        \end{align*}

  \item (2)

        证明与(1)类似。
\end{itemize}

\section*{6}
\begin{align*}
  \int_{a}^{a+p} f(x) dx
   & = \int_{a}^{0} f(x) dx + \int_{0}^{p} f(x) dx + \int_{p}^{a+p} f(x) dx   \\
   & = - \int_{0}^{a} f(x) dx + \int_{0}^{p} f(x) dx + \int_{p}^{a+p} f(x) dx
\end{align*}

令$x = t + p, dx = dt$，当$t$由$0$变为$a$时，$x$由$p$增到$a + p$，于是
\begin{align*}
  \int_{p}^{a+p} f(x) dx
   & = \int_{0}^{a} f(t + p) dt \\
   & = \int_{0}^{a} f(t) dt     \\
   & = \int_{0}^{a} f(x) dx
\end{align*}
所以
\begin{align*}
  \int_{a}^{a+p} f(x) dx
   & = - \int_{0}^{a} f(x) dx + \int_{0}^{p} f(x) dx + \int_{p}^{a+p} f(x) dx \\
   & = - \int_{0}^{a} f(x) dx + \int_{0}^{p} f(x) dx + \int_{0}^{a} f(x) dx   \\
   & = \int_{0}^{p} f(x) dx
\end{align*}

\section*{7}

\begin{itemize}
  \item (1)

        令$x = \frac{\pi}{2} - t, dx = -dt$，当$t$由$\frac{\pi}{2}$变为$0$时，$x$由$0$增到$\frac{\pi}{2}$，于是
        \begin{align*}
          \int_{0}^{\frac{\pi}{2}} f(sinx) dx
           & = - \int_{\frac{\pi}{2}}^{0} f(sin (\frac{\pi}{2} - t)) dt \\
           & = \int_{0}^{\frac{\pi}{2}} f(cos t) dt                     \\
           & = \int_{0}^{\frac{\pi}{2}} f(cos x) dx
        \end{align*}

  \item (2)

        令$x = \pi - t, dx = -dt$，当$t$由$\pi$变为$0$时，$x$由$0$增到$\pi$，于是
        \begin{align*}
          \int_{0}^{\pi} x f(sinx) dx
           & = - \int_{\pi}^{0} (\pi - t) f(sin(\pi - t)) dt                \\
           & = \int_{0}^{\pi}  (\pi - t) f(sin t) dt                        \\
           & = \pi \int_{0}^{\pi} f(sin t) dt - \int_{0}^{\pi} tf(sin t) dt \\
           & = \pi \int_{0}^{\pi} f(sin x) dx - \int_{0}^{\pi} xf(sin x) dx
        \end{align*}
        所以
        \begin{align*}
          2 \int_{0}^{\pi} x f(sinx) dx & = \int_{0}^{\pi} f(sin x) dx               \\
          \int_{0}^{\pi} x f(sinx) dx   & = \frac{\pi}{2} \int_{0}^{\pi} f(sin x) dx
        \end{align*}

\end{itemize}

\section*{8}
 (i)
\begin{align*}
  J(m,n) & = \int_{0}^{\frac{\pi}{2}} sin^m x cos^n x dx                                                                                                 \\
         & = \frac{1}{m + 1}\int_{0}^{\frac{\pi}{2}} cos^{n-1} x d(sin^{m + 1}x)                                                                         \\
         & = \frac{1}{m + 1}\left( cos^{n-1} x sin^{m + 1}x \bigg|_0^{\frac{\pi}{2}} + (n-1)\int_{0}^{\frac{\pi}{2}} sin^{m + 2}x cos^{n-2} x dx \right) \\
         & = \frac{n - 1}{m + 1} \int_{0}^{\frac{\pi}{2}} sin^{m + 2}x cos^{n-2} x dx                                                                    \\
         & = \frac{n - 1}{m + 1} \int_{0}^{\frac{\pi}{2}} sin^{m}x (1 - cos^2 x) cos^{n-2} x dx                                                          \\
         & = \frac{n - 1}{m + 1} \left(\int_{0}^{\frac{\pi}{2}} sin^{m}x cos^{n-2} x dx - \int_{0}^{\frac{\pi}{2}} sin^{m}x cos^{n} x dx\right)          \\
         & = \frac{n - 1}{m + 1} \left(J(m, n - 2) - J(m, n)\right)
\end{align*}
于是可得
\begin{align*}
  J(m,n) + \frac{n - 1}{m + 1} J(m, n) & = \frac{n - 1}{m + 1} J(m, n - 2) \\
  \frac{m + n}{m + 1} J(m, n)          & = \frac{n - 1}{m + 1} J(m, n - 2) \\
  J(m, n)                              & = \frac{n - 1}{m + n} J(m, n - 2)
\end{align*}
令$t = \frac{\pi}{2} - x, dx = -dt$，于是
\begin{align*}
  J(m, n) & = \int_{0}^{\frac{\pi}{2}} sin^m x cos^n x dx                                       \\
          & = - \int_{\frac{\pi}{2}}^{0} sin^m (\frac{\pi}{2} - t) cos^n (\frac{\pi}{2} - t) dt \\
          & = \int_{0}^{\frac{\pi}{2}} sin^m (\frac{\pi}{2} - t) cos^n (\frac{\pi}{2} - t) dt   \\
          & = \int_{0}^{\frac{\pi}{2}} cos^m t sin^n t dt                                       \\
          & = J(n, m)
\end{align*}
利用以上结论，我们有
\begin{align*}
  J(m, n) & = J(n, m)                         \\
          & = \frac{m - 1}{n + m} J(n, m - 2) \\
          & = \frac{m - 1}{m + n} J(m - 2, n)
\end{align*}

(ii)
\begin{align*}
  J(2m, 2n) & = \frac{2m - 1}{2m + 2n} J(2m - 2, 2n)                                                     \\
            & = \frac{2m - 1}{2m + 2n} \cdot \frac{2m - 3}{2m + 2n - 2} J(2m - 4, 2n)                    \\
            & = \frac{2m - 1}{2m + 2n} \cdot \frac{2m - 3}{2m + 2n - 2} \cdots \frac{1}{2n + 2} J(0, 2n) \\
            & = \frac{(2m - 1)!!}{(2m + 2n)(2m + 2n - 2)\cdot(2n + 2)} J(0, 2n)
\end{align*}
另外
\begin{align*}
  J(0, 2n) & = \frac{2n - 1}{2n} J(0, 2n - 2)                                           \\
           & = \frac{2n - 1}{2n} \cdot \frac{2n - 3}{2n - 2} J(0, 2n - 4)               \\
           & = \frac{2n - 1}{2n} \cdot \frac{2n - 3}{2n - 2} \cdots \frac{1}{2} J(0, 0) \\
\end{align*}
于是可得
\begin{align*}
  J(2m, 2n) & = \frac{(2m - 1)!!}{(2m + 2n)(2m + 2n - 2)\cdot(2n + 2)} J(0, 2n)                                                                 \\
            & = \frac{(2m - 1)!!}{(2m + 2n)(2m + 2n - 2)\cdot(2n + 2)} \frac{2n - 1}{2n} \cdot \frac{2n - 3}{2n - 2} \cdots \frac{1}{2} J(0, 0) \\
            & = \frac{(2m - 1)!!(2n - 1)!!}{(2m + 2n)!!}J(0, 0)                                                                                 \\
            & = \frac{(2m - 1)!!(2n - 1)!!}{(2m + 2n)!!} \cdot \frac{\pi}{2}
\end{align*}

\textbf{注: } 为了计算方便列出一下定积分的值
\begin{align*}
  J(0, 0) & = \int_{0}^{\frac{\pi}{2}} 1 dx = \frac{\pi}{2}        \\
  J(1, 0) & = \int_{0}^{\frac{\pi}{2}} sin x dx = 1                \\
  J(0, 1) & = \int_{0}^{\frac{\pi}{2}} cos x dx = 1                \\
  J(1, 1) & = \int_{0}^{\frac{\pi}{2}} sinx cos x dx = \frac{1}{2}
\end{align*}



\section*{9}

由习题1可知
\begin{align*}
  g'(x) = af(ax) - f(x) = 0
\end{align*}
即
\begin{align*}
  f(x) = a f(ax)
\end{align*}
对任意$x > 0$，取$a = \frac{1}{x}$，有
\begin{align*}
  f(x) & = \frac{1}{x} f(\frac{1}{x} x) \\
       & = \frac{1}{x} f(1)
\end{align*}
令$c = f(1)$，$f(x) = \frac{c}{x}$，命题成立。

\section*{10}

注意：题中的“连续可微函数”，指的是：导函数是连续的。

\section*{11}

设两个阴影从左到右分别为$S_1, S_2$，于是可得
\begin{align*}
  S_1 & = \int_{a}^{\xi} f(x) dx - \int_{a}^{\xi} f(a) dx = \int_{a}^{\xi} f(x) - f(a) dx \\
  S_2 & = \int_{\xi}^{b} f(b) dx - \int_{\xi}^{b} f(x) dx = \int_{\xi}^{b} f(b)- f(x) dx
\end{align*}
于是
\begin{align*}
  S_1 - S_2
   & =  \int_{a}^{\xi} f(x) - f(a) dx - \int_{\xi}^{b} f(b)- f(x) dx
\end{align*}
定义辅助函数
\begin{align*}
  F(x) = \int_{a}^{x} f(t) - f(a) dt - \int_{x}^{b} f(b)- f(t) dt
\end{align*}
由题设和微积分基本定理可知，$F(x)$是连续可导函数，且
\begin{align*}
  F(a) & = \int_{a}^{a} f(t) - f(a) dt - \int_{a}^{b} f(b)- f(t) dt \\
       & = - \int_{a}^{b} f(b)- f(t) dt
\end{align*}
由与$f(x)$是$[a, b]$上的严格增的连续函数，所以$f(b) - f(t) > 0, t \in [a, b)$， 可得
\begin{align*}
  F(a) & = - \int_{a}^{b} f(b)- f(t) dt < 0
\end{align*}
同理
\begin{align*}
  F(b) & = \int_{a}^{x} f(t) - f(a) dt - \int_{x}^{b} f(b)- f(t) dt  \\
       & = \int_{a}^{b} f(t) - f(a) dt - \int_{b}^{b} f(b) - f(t) dt \\
       & = \int_{a}^{b} f(t) - f(a) dt                               \\
       & > 0
\end{align*}
由连续函数的介质性定理可知，存在$\xi \in (a, b)$，使得
\begin{align*}
  \int_{a}^{\xi} f(x) - f(a) dx - \int_{\xi}^{b} f(b)- f(x) dx = 0
\end{align*}
即
\begin{align*}
  S_1 - S_2 = 0
\end{align*}
命题得证。

\section*{12}

\begin{itemize}
  \item 方法一：

        利用第二积分中值定理的推论，存在$\xi \in [0, 2\pi]$，使得
        \begin{align*}
          \int_{0}^{2\pi} f(x) sin nx dx
           & = f(0)\int_{0}^{\xi} sin nx dx + f(2\pi)\int_{\xi}^{2\pi} sin nx dx                               \\
           & = f(0) (- \frac{1}{n} cos nx\bigg|_{0}^{\xi}) + f(2\pi) (- \frac{1}{n} cos nx\bigg|_{\xi}^{2\pi}) \\
           & = f(0) [- \frac{1}{n} (cos n\xi - 1)] + f(2\pi) [-\frac{1}{n}(1 - cos n\xi)]                      \\
           & = (1 - cos n\xi) \frac{f(0)- f(2\pi)}{n}                                                          \\
           & \geq 0
        \end{align*}

  \item 方法二：

        令$h(x) = f(x) - f(2\pi)$，因为$f(x)$为单调递减函数，则$h$为非负、递减函数。由定理
        9.11(i)，存在$\xi \in [0, 2\pi]$，使得
        \begin{align*}
          \int_{0}^{2\pi} h(x) sin nx dx
           & = \int_{0}^{2\pi} (f(x) - f(2\pi)) sin nx dx \\
           & = h(0) \int_{0}^{\xi} sin nx dx
        \end{align*}
        因为
        \begin{align*}
          \int_{0}^{2\pi} f(x) sin nx dx
           & = \int_{0}^{2\pi} (f(x) - f(2\pi)) sin nx + f(2\pi)sin nx dx         \\
           & = \int_{0}^{2\pi} h(x) sin nx dx + f(2\pi) \int_{0}^{2\pi} sin nx dx \\
           & = h(0)\int_{0}^{\xi} sin nx dx + 0                                   \\
           & = h(0)(- \frac{1}{n} cos nx\bigg|_{0}^{\xi})                         \\
           & = h(0)(- \frac{1}{n} (cos n\xi - 1))                                 \\
           & = h(0)\frac{1 - cos n\xi}{n}                                         \\
           & \geq 0
        \end{align*}

\end{itemize}

\section*{13}

令$t^2 = u, t = \sqrt{u}, dt = \frac{1}{2\sqrt{u}}du$，于是
\begin{align*}
  \int_{x}^{x + c} sin t^2 dt
   & = \int_{x^2}^{(x + c)^2} \frac{sin u}{2\sqrt{u}}du
\end{align*}
由于$\frac{1}{2\sqrt{u}}$在$[x^2, (x + c)^2]$是单调递减的，且$\frac{1}{2\sqrt{u}} > 0$，
由积分第二中值定理可知，存在$\xi \in [x^2, (x + c)^2]$，使得
\begin{align*}
  \int_{x^2}^{(x + c)^2} \frac{sin u}{2\sqrt{u}}du
   & = \frac{1}{2\sqrt{x^2}} \int_{x^2}^{\xi} sin u du \\
   & = \frac{1}{2x} (- cos u)\bigg|_{x^2}^{\xi}        \\
   & = \frac{1}{2x} (cos x^2 - cos \xi)
\end{align*}
因为
\begin{align*}
  -2 \leq |cos x^2 - cos \xi| \leq 2
\end{align*}
所以
\begin{align*}
  - \frac{1}{x} \leq \left|\frac{1}{2x} (cos x^2 - cos \xi)\right| \leq \frac{1}{x}
\end{align*}
即
\begin{align*}
  \left| \int_{x}^{x + c} sin t^2 dt \right| \leq \frac{1}{x}
\end{align*}

\section*{14}
todo

首先要确定$f(\varphi(t))\varphi'(t)$是可积的。

\section*{15}

令
\begin{equation*}
  g'(x) =
  \begin{cases}
    f'(x), & f'(x) \geq 0 \\
    0,     & f'(x) < 0
  \end{cases},
  h'(x) =
  \begin{cases}
    f'(x), & f'(x) \leq 0 \\
    0,     & f'(x) > 0
  \end{cases},
\end{equation*}
因为$f$在$[a, b]$上是连续可微的，所以$g'(x), h'(x)$都是连续，
令$g(x), h(x)$分别是$\int_{a}^{x} g'(t) dt, \int_{a}^{x} h'(t) dt$一个原函数，
于是
\begin{align*}
  \int_{a}^{x} f'(t) dt & = \int_{a}^{x} g'(t) + h'(t) dt \\
  f(x) - f(a)           & = (g(t) + h(t))\bigg|_{a}^{x}   \\
                        & = g(x) + h(x) - g(a) - h(a)
\end{align*}
可得
\begin{align*}
  f(x) = g(x) + h(x) + f(a) - g(a) - h(a)
\end{align*}
取$C = f(a) - g(a) - h(a)$。
综上，$g(x)$是增函数，$h(x) = h(x) + C$任然是减函数。

\end{document}